% Bagian: computability
% Bab: computability-theory
% Subbagian: fixed-point-thm

\documentclass[../../../include/open-logic-section]{subfiles}

\begin{document}

\olfileid[id]{cmp}{thy}{fix}
\olsection{Teorema Titik Tetap}

Mari kita tinjau kembali masalah penghentian. Sebagai notasi sementara,
tulislah $\gn{\cfind{x}(y)}$ untuk $\tuple{x,y}$; bayangkan notasi ini sebagai
representasi suatu ``nama'' bagi nilai $\cfind{x}(y)$. Dengan notasi ini,
kita dapat merumuskan ulang salah satu bukti bahwa masalah penghentian tidak
dapat diputuskan.

Pertanyaan: adakah fungsi dapat dihitung $h$ dengan sifat berikut? Bagi setiap
$x$ dan $y$,
\[
h(\gn{\cfind{x}(y)}) =
\begin{cases}
1 & \text{jika $\cfind{x}(y) \fdefined$} \\
0 & \text{selain itu.}
\end{cases}
\]

Jawaban: Tidak; jika ada, fungsi parsial
\[
g(x) \simeq
\begin{cases}
0 & \text{jika $h(\gn{\cfind{x}(x)}) = 0$} \\
\text{tidak terdefinisi} & \text{selain itu}
\end{cases}
\]
akan dapat dihitung, sehingga mempunyai suatu indeks~$e$. Namun, kita lalu
mempunyai
\[
\cfind{e}(e) \simeq
\begin{cases}
0 & \text{jika $h(\gn{\cfind{e}(e)}) = 0$} \\
\text{tidak terdefinisi} & \text{selain itu,}
\end{cases}
\]
dan dalam keadaan ini $\cfind{e}(e)$ terdefinisi jika dan hanya jika tidak
terdefinisi, suatu kontradiksi.

Sekarang, perhatikan persamaan yang memuat $\cfind{e}$. Dalam suatu arti,
terdapat acuan diri di sana: kita telah mengatur agar nilai
$\cfind{e}(e)$ bergantung pada $\gn{\cfind{e}(e)}$ dengan cara tertentu.
Teorema titik tetap mengatakan bahwa secara umum kita \emph{dapat} melakukan
hal ini---bukan hanya demi membuktikan kontradiksi.

\olref{lem:fixed-equiv} memberikan dua cara ekuivalen untuk menyatakan teorema
titik tetap. Secara logis, fakta bahwa kedua pernyataan itu ekuivalen
mengikuti dari fakta bahwa keduanya benar; tetapi yang sebenarnya kita
maksudkan ialah bahwa masing-masing mengikuti secara langsung dari yang lain,
sehingga keduanya dapat dipandang sebagai pernyataan alternatif bagi teorema
yang sama.

\begin{lem}
\ollabel{lem:fixed-equiv}
Pernyataan-pernyataan berikut ekuivalen:
\begin{enumerate}
\item Bagi setiap fungsi dapat dihitung parsial $g(x,y)$, terdapat indeks~$e$
  sedemikian sehingga bagi setiap~$y$,
\[
\cfind{e}(y) \simeq g(e,y).
\]
\item Bagi setiap fungsi dapat dihitung~$f(x)$, terdapat indeks~$e$ sedemikian
  sehingga bagi setiap~$y$,
\[
\cfind{e}(y) \simeq \cfind{f(e)}(y).
\]
\end{enumerate}
\end{lem}

\begin{proof}
$(1) \Rightarrow (2)$: Diberikan $f$, definisikan $g$ dengan $g(x,y) \simeq
\fn{Un}(f(x),y)$. Gunakan (1) untuk memperoleh indeks~$e$ sedemikian sehingga
bagi setiap~$y$,
\begin{align*}
\cfind{e}(y) & \simeq \fn{Un}(f(e),y) \\
& \simeq \cfind{f(e)}(y).
\end{align*}

$(2) \Rightarrow (1)$: Diberikan $g$, gunakan teorema $s$-$m$-$n$ untuk
memperoleh $f$ sedemikian sehingga bagi setiap $x$ dan~$y$,
$\cfind{f(x)}(y) \simeq g(x,y)$. Gunakan (2) untuk memperoleh indeks~$e$
sedemikian sehingga
\begin{align*}
\cfind{e}(y) & \simeq \cfind{f(e)}(y) \\
& \simeq g(e,y).
\end{align*}
Ini menyelesaikan bukti.
\end{proof}

\begin{explain}
Sebelum menunjukkan bahwa pernyataan (1) benar (dan karena itu (2) juga),
tinjau betapa ganjilnya pernyataan itu. Bayangkan $e$ sebagai program
komputer; pernyataan (1) mengatakan bahwa, diberikan sembarang fungsi dapat
dihitung parsial $g(x,y)$, Anda dapat menemukan program komputer~$e$ yang
menghitung $g_e(y) \simeq g(e,y)$. Dengan kata lain, Anda dapat menemukan
program komputer yang menghitung fungsi yang mengacu pada program itu sendiri.
\end{explain}

\begin{thm}
Kedua pernyataan dalam \olref{lem:fixed-equiv} benar. Secara khusus, bagi
setiap fungsi dapat dihitung parsial $g(x,y)$, terdapat indeks~$e$ sedemikian
sehingga bagi setiap~$y$,
\[
\cfind{e}(y) \simeq g(e,y).
\]
\end{thm}

\begin{proof}
Bahan-bahannya sudah tersirat dalam pembahasan masalah penghentian di atas.
Misalkan $\fn{diag}(x)$ fungsi dapat dihitung yang bagi setiap $x$
mengembalikan indeks bagi fungsi $f_x(y) \simeq \cfind{x}(x,y)$, yakni
\[
\cfind{\fn{diag}(x)}(y) \simeq \cfind{x}(x,y).
\]
Bayangkan $\fn{diag}$ sebagai fungsi yang mengubah program bagi fungsi
berargumen~2 menjadi program bagi fungsi berargumen~1, yang diperoleh dengan
menetapkan program semula sebagai argumen pertamanya. Fungsi $\fn{diag}$
dapat didefinisikan secara formal sebagai berikut: mula-mula definisikan $s$
dengan
\[
s(x,y) \simeq \fn{Un}^2(x,x,y),
\]
dengan $\fn{Un}^2$ fungsi berargumen~3 yang universal bagi fungsi-fungsi
dapat dihitung parsial berargumen~2. Kemudian, menurut teorema $s$-$m$-$n$,
kita dapat menemukan fungsi rekursif primitif $\fn{diag}$ yang memenuhi
\[
\cfind{\fn{diag}(x)}(y) \simeq s(x,y).
\]

Sekarang, definisikan fungsi $l$ dengan
\[
l(x,y) \simeq g(\fn{diag}(x),y).
\]
Misalkan $\gn{l}$ suatu indeks bagi~$l$. Terakhir, misalkan
$e=\fn{diag}(\gn{l})$. Maka bagi setiap $y$, kita mempunyai
\begin{align*}
\cfind{e}(y) & \simeq \cfind{\fn{diag}(\gn{l})}(y) \\
& \simeq \cfind{\gn{l}}(\gn{l}, y) \\
& \simeq l(\gn{l}, y) \\
& \simeq g(\fn{diag}(\gn{l}),y) \\
& \simeq g(e, y),
\end{align*}
sebagaimana diperlukan.
\end{proof}

\begin{explain}
Apa yang sedang terjadi? Andaikan Anda diberi tugas menulis program komputer
yang mencetak dirinya sendiri. Namun, andaikan pula bahwa Anda bekerja dengan
bahasa pemrograman yang mempunyai pustaka fungsi untai yang kaya dan ganjil.
Secara khusus, andaikan bahasa pemrograman Anda mempunyai fungsi $\fn{diag}$
yang bekerja sebagai berikut: diberikan untai masukan~$s$, $\fn{diag}$
menemukan setiap kemunculan simbol `x' dalam~$s$, lalu menggantinya dengan
versi berpetik dari untai semula. Sebagai contoh, jika diberikan untai
\begin{quote}
\begin{verbatim}
hello x world
\end{verbatim}
\end{quote}
sebagai masukan, fungsi itu mengembalikan
\begin{quote}
\begin{verbatim}
hello 'hello x world' world
\end{verbatim}
\end{quote}
sebagai keluaran. Dalam hal ini, program yang diinginkan mudah ditulis; Anda
dapat memeriksa bahwa
\begin{quote}
\begin{verbatim}
print(diag('print(diag(x))'))
\end{verbatim}
\end{quote}
berhasil melakukannya. Untuk bahasa pemrograman yang lebih umum seperti C++
dan Java, gagasan yang sama tetap bekerja (dengan implementasi yang lebih
rumit).

Kita hanya tinggal beberapa langkah dari bukti teorema titik tetap. Andaikan
suatu varian fungsi cetak, $\fn{print}(x,y)$, menerima untai~$x$ dan argumen
numerik lain~$y$, lalu mencetak untai~$x$ sebanyak~$y$ kali.
Maka ``program''
\begin{quote}
\begin{verbatim}
getinput(y);
print(diag('getinput(y); print(diag(x), y)'), y)
\end{verbatim}
\end{quote}
mencetak dirinya sendiri sebanyak~$y$ kali ketika menerima masukan~$y$.
Mengganti kerangka $\fn{getinput}$---$\fn{print}$---$\fn{diag}$ dengan fungsi
sembarang $g(x,y)$ menghasilkan
\begin{quote}
\begin{verbatim}
g(diag('g(diag(x), y)'), y)
\end{verbatim}
\end{quote}
yang merupakan program yang, ketika menerima masukan~$y$, menjalankan $g$ pada
program itu sendiri dan~$y$. Jika kita memahami ``pengutipan'' sebagai
``penggunaan suatu indeks bagi,'' kita memperoleh bukti di atas.

Untuk sekarang, tidak apa-apa jika Anda ingin memandang bukti tersebut
sebagai siasat formal atau ilmu gaib. Namun, Anda harus mampu merekonstruksi
rincian argumen yang diberikan di atas. Ketika kita membuktikan teorema
ketaklengkapan (dan ``teorema titik tetap'' terkait), kita akan membahas
cara lain untuk memahami mengapa argumen ini berhasil.
\end{explain}

\begin{tagblock}{lambda}
\begin{digress}
Gagasan yang sama dapat digunakan untuk memperoleh kombinator ``titik
tetap.'' Andaikan Anda mempunyai suku lambda~$g$ dan menginginkan suku lain
$k$ dengan sifat bahwa $k$ ekuivalen-$\beta$ dengan $gk$. Definisikan suku
\[
\fn{diag}(x) = xx
\]
dan
\[
l(x) = g(\fn{diag}(x))
\]
menurut konvensi notasi kita; dengan kata lain, $l$ adalah suku
$\lambd[x][g(xx)]$. Misalkan $k$ suku~$ll$. Maka kita mempunyai
\begin{align*}
k & = (\lambd[x][g(xx)])(\lambd[x][g(xx)]) \\
& \red  g((\lambd[x][g(xx)])(\lambd[x][g(xx)])) \\
& = gk.
\end{align*}
Jika kita mengambil
\[
Y = \lambd[g][((\lambd[x][g(xx)])(\lambd[x][g(xx)]))]
\]
maka $Yg$ dan $g(Yg)$ tereduksi ke suatu suku yang sama; jadi
$Yg \equiv_\beta g(Yg)$. Ini dikenal sebagai ``kombinator Curry.'' Jika
sebagai gantinya kita mengambil
\[
Y = (\lambd[xg][g(xxg)])(\lambd[xg][g(xxg)])
\]
maka sebenarnya $Yg$ tereduksi ke $g(Yg)$, yang merupakan pernyataan lebih
kuat. Versi $Y$ yang terakhir ini dikenal sebagai ``kombinator Turing.''
\end{digress}
\end{tagblock}

\end{document}
